\begin{table}[htbp]
\centering
\footnotesize
\setlength{\tabcolsep}{4pt}
\caption{本研究各分析轮次采用的设定规格，以及 2026-10-20 数据版本下的 95\% 置信区间。}
\label{tab:s9}
\begin{tabular}{p{0.32\linewidth}p{0.24\linewidth}ll}
\toprule
轮次与规格 & 定位 & $a$ [95\% CI] & $b/a$ [95\% CI] \\
\midrule
第 1 轮：$G$，三个日历窗口（$S=3$） & 第 1 轮主要规格；第 2 轮起为次要规格 & 0.0991 [0.0569, 0.1382] & 22.6 [9.0, 52.2] \\
第 2 轮，规则第 1 版：在流行峰值结束窗口 & 废弃（2024--25 季仅四周）；未保留估计值 & -- & -- \\
第 2 轮，规则第 2 版（2026-10-20）：$G$，规则窗口（$S=4$） & 次要规格 & 0.0840 [0.0498, 0.1248] & 23.5 [8.4, 59.7] \\
第 2 轮，规则第 2 版：$R_A$，规则窗口（$S=4$） & 第 2 轮起的主要规格 & 0.0150 [0.0100, 0.0208] & 37.4 [19.1, 71.0] \\
第 2 轮：$R_A$，日历窗口（$S=3$） & 敏感性检验 & 0.0202 [0.0097, 0.0306] & 25.6 [10.2, 79.6] \\
第 3 轮：两步模拟反演 & 第 2 轮拟合的校准；无重新拟合 & -- & -- \\
\bottomrule
\end{tabular}
\par\smallskip
\begin{minipage}{0.95\linewidth}\scriptsize 全部估计值均在同一数据版本上以双向自助法重新拟合（$B=2000$）。更新型比值在第 2 轮中成为主要统计量，因其通过代际间隔与 $R_t$ 建立显式联系；加权增长比为保持连贯性而予以保留。各规格的点估计均落在同一统计量另一规格的置信区间之内。\end{minipage}
\end{table}
